\documentclass[../Main.tex]{subfiles}

\begin{document}
\section{Concepts and Definitions}
\begin{definition}{State}
    A quantum system is described at each instant by a \underline{state} $\state$ which belongs to a complex vector space $V$. Any non-zero element of this space $V$ is a possible state.
\end{definition}
Since these are in a vector space, for $\state, \ket{\phi} \in V$, $\alpha \state + \beta \ket{\phi} \in V$ for any complex constants $\alpha, \beta$.

Physically, this is a superposition principle leading to wave-like behaviour, but these states are not wavefunctions. There are also dual or conjugate states $\bra{\psi}$ which belong to the dual space $V^\dagger$. By definition, states and duals can be combined to give a complex number:
\begin{equation*}
    \bra{\phi} \text{ combined with }\state \mapsto \inn{\phi}{\psi}
\end{equation*}
The spaces $V$ and $V^\dagger$ are equipped with an inner product which can be described as a one-to-one correspondence between the spaces:
\begin{align*}
    V &\mapsto V^\dagger \\
    \state &\mapsto \bra{\psi} = \left(\state\right)^\dagger
\end{align*}
Then the inner product requires the first state to be mapped to its dual, then applied to the second:
\begin{equation*}
    \ket{\phi},~~\state \mapsto \inn\phi\psi = \left(\ket{\phi}\right)^T \state \in \C
\end{equation*}

\begin{remarks}
    \item Knowing $\inn\phi\psi$ for all $\bra\phi$ determines $\ket\psi$. \label{remInnDefinesState}
    \item The physical content of any state is unchanged by $\state \mapsto \alpha \state$ for $\alpha \in \C \backslash \{0\}$.
    \item We will usually normalise the state to make $\norm\state = 1$, but this still allows for relative phases $\state \mapsto e^{i\theta}\state$, which can be important for superposition.
    \item The space $V$ is complete. We will assume that appropriate sequences and series converge.
    \item The space of states is also known as a Hilbert space $\hilb$.
\end{remarks}
\begin{definition}{Operator}
    An \underline{operator} $Q$ is a linear map on states:
    \begin{align*}
        Q : V &\mapsto V\\
        \state &\mapsto Q\state
    \end{align*}
    which must obey linearity.
\end{definition}
The same operator can also be thought of acting to the left on dual states:
\begin{align*}
    Q : V^\dagger &\mapsto V^\dagger\\
    \bra{\psi} &\mapsto \bra{\psi}Q
\end{align*}
where the RHS is defined by remark~\ref{remInnDefinesState} as:
\begin{equation*}
    \left(\ket{\phi}Q\right)\state = \ket{\phi} \left(Q\state\right) = \opbraket{\phi}{Q}{\psi}
\end{equation*}
\end{document}